%! Absolute Value Equations & Inequalities — Unit 1, Lesson 1.2 — Homework (Answer Key)
% Mirrors homework/main.tex exactly; every \blank{} becomes an \ans{}, every \writelines{n} is
% answered with exactly n \ansline{}s, and every work block is authored byte-identically.
\documentclass[12pt]{article}
\usepackage{algebra2-article}
\usepackage{algebra2-key}   % pulls in -boxes; do NOT also load -boxes
\usepackage{grading}
\gradingassignment{alg2-u01-l02-hw}
\gradingcourse{algebra_2}
\usepackage{algebra2-graded}
% Plain number line: {xmin}{xmax}, ticks labelled every unit.
\newcommand{\numline}[2]{%
  \draw[->, semithick, charcoal] (#1-0.5,0)--(#2+0.5,0);
  \foreach \x in {#1,...,#2}{\draw[charcoal] (\x,0.12)--(\x,-0.12) node[below, font=\scriptsize]{$\x$};}}

\begin{document}

\pageheader{Unit 1, Lesson 1.2}{Homework --- Answer Key}
% NAMESTRIP: no \namedateperiod here — the name row lives on the cover only.

\vspace{0.10in}

\boxguard[8]
\begin{remindbox}
\small
\textbf{This is your graded homework.} Your packet with completed homework is due the first
class after two study halls. I will announce the due date in class and on TurtleNet. Remember: $|x-h|$ is a \emph{distance}, so \textbf{isolate
the bars before you split}, ``less th\emph{AND}'' gives one interval \emph{between}, and
``great\emph{OR}'' gives two rays \emph{outside}.
\end{remindbox}

\vspace{0.10in}

\boxguard[14]
\begin{notesbox}{Practice}
\small
\begin{enumerate}[leftmargin=1.9em, itemsep=7pt]
  \item \textbf{Solve two equations.} The first splits immediately; the second does not.
    \begin{enumerate}[label=(\alph*), leftmargin=1.9em, itemsep=3pt]
      \item $|x - 7| = 4$: \quad $x =$ \inblank{\blank{3.0cm}}{\gradedblank[id=q1a, type=short, points=1, answer={11, 3}, check=set, display={$11$ or $3$}]{3.0cm-2pt}{12.5pt}}
      \item $5|x + 2| - 1 = 14$. Isolate first.
        \begin{work}
          5|x+2| - 1 &= 14 \\
          5|x+2| &= 15 \\
          |x+2| &= 3 \\
          x + 2 &= 3 \quad\text{or}\quad x + 2 = -3 \\
          x &= 1 \quad\text{or}\quad x = -5
        \end{work}
        $x =$ \inblank{\blank{3.0cm}}{\gradedblank[id=q1b, type=short, points=2, answer={1, -5}, check=set, display={$1$ or $-5$}]{3.0cm-2pt}{12.5pt}}
    \end{enumerate}

  \boxguard[16]
  \item \textbf{One symbol apart.} Solve both, and notice what changes.
    \begin{center}
    \renewcommand{\arraystretch}{1.4}
    \begin{tabularx}{\linewidth}{@{}X|X@{}}
    \textbf{(a) $|x + 3| < 4$} & \textbf{(b) $|x + 3| > 4$} \\
    \hline
    $-4 < x+3 < 4$, so \inblank{\blank{2.4cm}}{\gradedblank[id=q2a, type=short, points=1, answer={-7<x<1}, check=expression, display={$-7 < x < 1$}]{2.4cm-2pt}{12.5pt}} & $x+3 < -4$ or $x+3 > 4$, so \inblank{\blank{2.6cm}}{\gradedblank[id=q2b, type=short, points=1, answer={x<-7 or x>1}, check=expression, display={$x<-7$ or $x>1$}]{2.6cm-2pt}{12.5pt}} \\
    Interval: \inblank{\blank{2.2cm}}{\gradedblank[id=q2c, type=short, points=1, answer={(-7,1)}, check=text]{2.2cm-2pt}{12.5pt}} & Interval: \inblank{\blank{3.4cm}}{\gradedblank[id=q2d, type=short, points=1, answer={(-inf,-7)U(1,inf)}, check=text, display={$(-\infty,-7)\cup(1,\infty)$}]{3.4cm-2pt}{12.5pt}} \\
    \end{tabularx}
    \end{center}
    (c) One of these is a single connected piece and the other is two. Say which, and why the
    \emph{distance} meaning makes that happen.
    \gradedlines{2}{\gradedblank[id=q2e, type=rubric, points=2, display={(a) is one piece: ``closer than 4 to $-3$'' is everything between the endpoints. (b) is two: ``farther than 4'' can overshoot on either side, so it splits.}]{\linewidth}{\gradedlinesheight}}

  \boxguard[16]
  \item \textbf{Read it off the picture.} The shaded numbers below are the solution set of an
        absolute value inequality.
    \begin{center}
    \begin{tikzpicture}[x=0.50cm, y=0.46cm]
      \numline{-4}{10}
      \draw[very thick, forest] (-1,0)--(7,0);
      \fill[forest] (-1,0) circle (3.0pt);
      \fill[forest] (7,0) circle (3.0pt);
    \end{tikzpicture}
    \end{center}
    \vspace{-4pt}
    (a) The centre of the shading is at \inblank{\blank{1.2cm}}{\gradedblank[id=q3a, type=short, points=1, answer={3}, check=numeric]{1.2cm-2pt}{12.5pt}}, and it reaches \inblank{\blank{1.2cm}}{\gradedblank[id=q3b, type=short, points=1, answer={4}, check=numeric]{1.2cm-2pt}{12.5pt}} units each
    way. \\[4pt]
    (b) So the inequality is \inblank{\blank{3.0cm}}{\gradedblank[id=q3c, type=short, points=1, answer={|x-3|<=4}, check=expression, display={$|x-3| \le 4$}]{3.0cm-2pt}{12.5pt}}, and in interval notation the set is
    \inblank{\blank{2.4cm}}{\gradedblank[id=q3d, type=short, points=1, answer={[-1,7]}, check=text]{2.4cm-2pt}{12.5pt}}

  \boxguard[16]
  \item \textbf{The odd ones out.} Answer each from the \emph{distance} meaning, not a memorised
        rule.
    \begin{enumerate}[label=(\alph*), leftmargin=1.9em, itemsep=3pt]
      \item $|x - 4| = -3$ has \inblank{\blank{2.8cm}}{\gradedblank[id=q4a, type=short, points=1, answer={no solution}, check=text]{2.8cm-2pt}{12.5pt}} \quad $|x - 4| = 0$ has \inblank{\blank{2.8cm}}{\gradedblank[id=q4b, type=short, points=1, answer={one}, check=text, display={one, $x=4$}]{2.8cm-2pt}{12.5pt}}
      \item $|x - 4| \ge -3$ has \inblank{\blank{3.2cm}}{\gradedblank[id=q4c, type=short, points=1, answer={all real numbers}, check=text]{3.2cm-2pt}{12.5pt}} \; because \inblank{\blank{4.0cm}}{\gradedblank[id=q4d, type=rubric, points=2, display={distance $\ge 0 > -3$ always}]{4.0cm-2pt}{12.5pt}}
      \item For which values of $k$ does $|x - 3| < k$ have \emph{no} solution? \inblank{\blank{2.0cm}}{\gradedblank[id=q4e, type=short, points=1, answer={k<=0}, check=expression, display={$k \le 0$}]{2.0cm-2pt}{12.5pt}}
            \; Be careful at the boundary.
            \rlap{\hskip1pt\raisebox{0pt}[0pt][0pt]{\gradedblank[id=q4f, type=rubric, points=1, display={At $k=0$ it reads ``distance $<0$'', still impossible --- so $k=0$ is included.}]{383.75pt}{12.5pt}}}\writelines{1}
    \end{enumerate}

  \boxguard[18]
  \item \textbf{Model a tolerance.} A thermostat holds a room within $2^\circ$F of its $68^\circ$F
        setting. The desired value is \inblank{\blank{1.4cm}}{\gradedblank[id=q5a, type=short, points=1, answer={68}, check=numeric, display={$68^\circ$F}]{1.4cm-2pt}{12.5pt}} and the tolerance is \inblank{\blank{1.4cm}}{\gradedblank[id=q5b, type=short, points=1, answer={2}, check=numeric, display={$2^\circ$F}]{1.4cm-2pt}{12.5pt}}.
    \begin{enumerate}[label=(\alph*), leftmargin=1.9em, itemsep=3pt]
      \item Write an absolute value inequality for an acceptable temperature $T$. \inblank{\blank{3.6cm}}{\gradedblank[id=q5c, type=short, points=1, answer={|T-68|<=2}, check=expression, display={$|T-68| \le 2$}]{3.6cm-2pt}{12.5pt}}
      \item Solve it.
        \begin{work}
          |T - 68| &\le 2 \\
          -2 \le T - 68 &\le 2 \\
          66 \le T &\le 70
        \end{work}
        Coolest \inblank{\blank{1.6cm}}{\gradedblank[id=q5d, type=short, points=1, answer={66}, check=numeric, display={$66^\circ$}]{1.6cm-2pt}{12.5pt}} \qquad Warmest \inblank{\blank{1.6cm}}{\gradedblank[id=q5e, type=short, points=1, answer={70}, check=numeric, display={$70^\circ$}]{1.6cm-2pt}{12.5pt}}
      \item Is $66^\circ$F acceptable? \inblank{\blank{1.4cm}}{\gradedblank[id=q5f, type=short, points=1, answer={yes}, check=text]{1.4cm-2pt}{12.5pt}} \; Which word in the sentence tells you?
            \inblank{\blank{2.4cm}}{\gradedblank[id=q5g, type=short, points=1, answer={within}, check=text, display={``within''}]{2.4cm-2pt}{12.5pt}}
    \end{enumerate}

  \boxguard[14]
  \item \textbf{Multiple choice (SOL-style).} Which is the solution set of $|x + 2| > 3$?
    \begin{enumerate}[label=\Alph*., leftmargin=2.2em, itemsep=1pt]
      \item \letterbox{\gradedblank[id=q6a, type=choice, points=1, answer={B}]{0.4in}{12.5pt}}$-5 < x < 1$
      \item $x < -5$ \; or \; $x > 1$ \hfill \textcolor{keyred}{\textbf{$\leftarrow$ correct}}
      \item $x > 1$
      \item all real numbers
    \end{enumerate}
    Say in one sentence how you ruled out one of the others.
    \rlap{\hskip1pt\raisebox{0pt}[0pt][0pt]{\gradedblank[id=q6b, type=rubric, points=2, display={A answers $|x+2|<3$ --- ``greater than'' points outward, not inward.}]{181.25pt}{12.5pt}}}\writelines{1}
\end{enumerate}
\end{notesbox}

\vspace{0.10in}

\boxguard[10]
\begin{spiralbox}
\small
\textbf{Coming next --- Lesson 1.3: Absolute Value Functions \& Transformations.} Today you
\emph{solved} $|ax+b| = c$; next you will \emph{graph} $y = |ax+b|$ --- a \textbf{V-shape}. The two
solutions you found for an equation are exactly where that V crosses the horizontal line $y = c$.
\end{spiralbox}

\end{document}
